\chapter{Conclusão}\label{cap:conclusao}

\textit{Capítulo a escrever depois dos resultados finais dos Capítulos~\ref{cap:micro} e~\ref{cap:macro}.}

Achados preliminares a retomar:
\begin{itemize}
  \item O saldo de empregos em Alojamento e Alimentação no Pará ficou abaixo do normal após a COP30 em todas
        as especificações da regressão com dummies de fase (Capítulo~\ref{cap:macro}), com significância apenas na
        Região Metropolitana de Belém. O painel de capitais (Capítulo~\ref{cap:micro}) testa se esse padrão se mantém.
  \item No mês do evento, a hotelaria teve aumento significativo de contratações.
  \item O padrão observado é de desaceleração, e não de destruição líquida de vagas: o saldo acumulado de
        dezembro a julho caiu, mas permaneceu positivo.
\end{itemize}

Também a retomar: somado de outubro a dezembro, o saldo de 2025 ($+957$) ficou acima de 2023 ($+126$) e
2024 ($+547$), o que indica que o aquecimento antes do evento não é apenas um deslocamento de calendário
provocado pelo Círio.

Pontos a discutir: o legado da infraestrutura, a informalidade que o CAGED não observa, os limites de um
efeito estimado com poucos meses pós-evento e as implicações para a decisão de sediar grandes eventos.
